\begin{textsample}{NEG Dim 9 – global\_south\_2024\_11 – Score -2.00 – global\_south\_2024\_11/117541288.txt}  \label{ex:f9_neg_309}
Saudi Crown Prince Mohammed bin Salman has condemned Israel ' s actions in Gaza as"genocide"in some of the harshest public criticism of the country by a Saudi official since the start of the war.
Speaking at a summit of Muslim and Arab leaders the prince also criticised Israeli attacks on Lebanon and Iran, and - in a sign of improving ties between rivals Riyadh and Tehran, warned that Israel should not launch attacks on Iranian soil.
Saudi ' s de facto leader was joined by other leaders present in calling for a total Israeli withdrawal from the West Bank and Gaza.
Meanwhile, Saudi Arabia ' s foreign minister said it was a"failing of the international community"that the war in Gaza had not been stopped, accusing Israel of causing starvation there.
Prince Faisal Bin Farhan Al-Saud said :"Where the international community primarily has failed is ending the immediate conflict and putting an end to Israel ' s aggression."The war in Gaza was triggered by Hamas ' s 7 October @ @ @ @ @ @ @ @ @ @ Israel.
About 1,200 people were killed and 251 others were taken hostage.
Israel retaliated by launching a military campaign to destroy Hamas, during which more than 43,400 people have been killed in Gaza, according to the Hamas-run health ministry.
Leaders at the summit also condemned what they described as Israel ' s"continuous attacks"against UN staff and facilities in Gaza.
Last month, the Knesset passed a bill to ban the Unrwa, the UN Palestinian refugee agency, from operating in Israel and occupied East Jerusalem, accusing the organisation of colluding with Hamas.
Several countries, including the US and the UK, have expressed serious concern about the move limiting the agency ' s ability to transfer aid to Gaza.
In the backdrop of the well-attended summit, is Donald Trump ' s return to the White House.
Gulf leaders are aware of his closeness to Israel, but they also have good relations with him, and want him to use his influence and his fondness for deal-making to secure an @ @ @ @ @ @ @ @ @ @ Arabia, Trump is viewed much more favourably than Joe Biden, but his track record in the Middle East is mixed.
He pleased Israel and angered the Muslim world by recognising Jerusalem as Israel ' s capital as well as the annexation of the occupied Golan Heights.
He also secured the Abraham Accords in 2020 which saw the UAE, Bahrain and Morocco establish full diplomatic relations with Israel and Sudan agree to do so.
Still, one editorial in a leading Saudi newspaper today is titled :"A new era of hope.
Trump ' s return and the promise of stability."DISCLAIMER : The Views, Comments, Opinions, Contributions and Statements made by \textbf{Readers} and Contributors on this platform do not \textbf{necessarily} represent the views or policy of Multimedia Group Limited.

% matched lemmas: necessarily, reader
\end{textsample}
